\begin{abstract}
We introduce \textsc{QLoRA}, a memory-efficient finetuning method that enables the adaptation of a 65B parameter language model using a single 48GB GPU, all while retaining the task performance typically associated with full 16-bit finetuning. \textsc{QLoRA} works by passing gradients through a frozen, 4-bit quantized backbone pretrained model into Low Rank Adapters (LoRA). Our top-performing model family, dubbed \textbf{Guanaco}, surpasses all prior publicly released models on the Vicuna benchmark, achieving 99.3\% of ChatGPT’s performance with only 24 hours of single-GPU finetuning. \textsc{QLoRA} advances memory savings without degrading effectiveness by introducing (a) 4-bit NormalFloat (NF4), an information-theoretically optimal data format for normally distributed parameters, (b) Double Quantization, which further lowers memory needs by also quantizing quantization parameters, and (c) Paged Optimizers, designed to alleviate peak memory usage. With \textsc{QLoRA}, we finetuned over 1,000 models, delivering an extensive assessment of instruction following and chatbot capabilities across 8 instruction datasets, several architectures (LLaMA, T5), and large parameter counts previously impractical for conventional finetuning (e.g., models with 33B and 65B parameters). Our experiments demonstrate that \textsc{QLoRA} achieves state-of-the-art performance on small, high-quality datasets, outperforming the former best results even when applied to relatively compact models. We offer an in-depth examination of chatbot efficacy, employing both human ratings and GPT-4-based assessments—finding the latter to be a cost-effective and valid proxy for human evaluation. Additionally, we observe that prevailing chatbot benchmarks do not reliably reflect actual chatbot capabilities. Through a lemon-picked evaluation, we highlight scenarios in which \textbf{Guanaco} underperforms compared to ChatGPT. All models, source code, and 4-bit training CUDA kernels are openly released.\footnote{\url{https://github.com/artidoro/qlora} and \url{https://github.com/TimDettmers/bitsandbytes}}
\end{abstract}

\section{Introduction}

Customizing large language models (LLMs) through finetuning is widely recognized as a potent method to enhance model performance \citep{min2021metaicl, wei2021finetuned, ouyang2022training, wang2022super, wang2022self, liu2022few} and to impart or suppress specific behaviors \citep{ouyang2022training,askell2021general,bai2022training}. Nevertheless, finetuning extremely large models presents substantial computational and financial challenges; for example, conducting standard 16-bit finetuning on a LLaMA model with 65 billion parameters~\citep{touvron2023llama} demands over 780 GB of GPU memory. Although recent advancements in quantization approaches have succeeded in lowering the memory consumption for LLM inference \citep{dettmers2022llm, dettmers2022case, frantar2022gptq, xiao2022smoothquant}, these strategies are restricted to inference and generally fail when applied during the training phase \citep{wortsman2023stable}.

In this work, we establish, for the first time, that finetuning a 4-bit quantized model can be achieved without any loss in model quality. The proposed method, \textsc{QLoRA}, introduces an innovative high-precision quantization technique for pretrained models to 4-bit precision, followed by the incorporation of a compact set of trainable Low-rank Adapter weights \citep{hu2021lora}.

These adapter weights are updated through gradient backpropagation through the quantized model parameters.

\textsc{QLoRA} brings the memory overhead required to finetune a 65B parameter model down from over $780$GB to less than $48$GB of GPU memory, all without compromising on training efficiency or accuracy relative to a 16-bit full finetuning setup. This development represents a substantial leap in the accessibility of LLM finetuning, rendering it possible to finetune the largest open models on a single GPU. Leveraging \textsc{QLoRA}, we developed the \textbf{Guanaco} series of models; notably, the second-best model attains 97.8\% of ChatGPT's performance as measured on the Vicuna benchmark~\citep{vicuna2023}, with training requiring under 12 hours on a consumer-grade GPU. Utilizing one professional-grade GPU for 24 hours, our most substantial model achieves 99.3\%, effectively matching ChatGPT’s performance on the same benchmark. Furthermore, when deployed, the smallest \textbf{Guanaco} model, featuring 7 billion parameters, operates with a memory footprint of just 5 GB and surpasses the 26 GB Alpaca model by more than 20 percentage points on the Vicuna benchmark (refer to Table~\ref{tbl:vicuna_eval}).

\textsc{QLoRA} presents several advancements aimed at minimizing memory consumption while maintaining high performance levels: (1) \textbf{4-bit NormalFloat}, an information-theoretically optimal quantization type tailored for data following a normal distribution, surpasses 4-bit Integers and 4-bit Floats in empirical evaluations. (2) \textbf{Double Quantization}, a technique that further compresses the quantization coefficients themselves, leading to an average reduction of 0.37 bits per parameter (equivalent to approximately 3 GB for models with 65B parameters). (3) \textbf{Paged Optimizers}, which employ NVIDIA unified memory to eliminate memory surges from gradient checkpointing during processing of mini-batches with extended sequence lengths. Integrating these components results in an improved LoRA methodology, placing adapters at all layers of the network, thereby alleviating nearly all the accuracy compromises encountered in earlier research.

Thanks to its memory efficiency, \textsc{QLoRA} enables extensive experimentation with instruction finetuning and the assessment of chatbot performance at scales previously unattainable with conventional finetuning methods due to their memory requirements. Consequently, we train over 1,000 models, spanning diverse instruction tuning corpora, architectures, and parameter counts ranging from 80M up to 65B. Our work demonstrates that \textsc{QLoRA} can match 16-bit baseline performance (\S\ref{sec:comp_to_std}) and is capable of training state-of-the-art chatbots such as \textbf{Guanaco} (\S\ref{sec:pushing}), while also enabling an analysis of emergent patterns within these models. Notably, our findings indicate that the quality of training data is significantly more critical than the sheer size of the dataset; for instance, a 9k-example dataset (OASST1) surpassed the performance of a subsampled 450k-example dataset (FLAN v2) in chatbot tasks, even though both were intended for instruction-following generalization. Additionally, we observe that strong results on the Massive Multitask Language Understanding (MMLU) benchmark do not necessarily translate to superior performance on the Vicuna chatbot benchmark and vice versa, highlighting that the appropriateness of the dataset for the task is more pivotal than its volume.

In addition, we conduct a comprehensive evaluation of chatbot systems employing both human evaluators and GPT-4 to assess performance. Our methodology leverages a tournament-style benchmarking framework, in which models are pitted against each other in direct competition, each attempting to generate the superior response to a given prompt. Outcomes are determined either by the judgment of GPT-4 or by human raters, and these results are subsequently synthesized into Elo ratings~\citep{elo1967proposed, elo1978rating}, which establish a hierarchy of chatbot effectiveness. Our findings reveal a high degree of concordance between GPT-4 and human judgments regarding overall model rankings; nevertheless, significant discrepancies do occur in some cases. This underlines that while model-based assessments can serve as cost-efficient proxies for human evaluation, they are also subject to distinct sources of uncertainty.

Beyond quantitative benchmarking, we complement our evaluation with a qualitative review of the \textbf{Guanaco} models. Through this qualitative lens, we identify both exemplary and deficient behaviors that escape detection by purely numerical scores.

To support further research, we provide public access to all model outputs along with corresponding annotations by both humans and GPT-4. Our entire codebase and CUDA kernel implementations are open-sourced and incorporated within the Hugging Face transformers library~\citep{wolf2019huggingface}, simplifying community access. We also publish a suite of adapters compatible with 7/13/33/65B model sizes, each trained using eight distinct instruction-following datasets, yielding 32 fully open, finetuned model variants.

\section{Background}
\paragraph{Block-wise k-bit Quantization} Quantization refers to the transformation of data from a format with higher information density to one with lower information density, commonly entailing a reduction in bit-width—for instance, converting from 32-bit floating-point values to 8-bit integers. To maximize utilization of the available numeric range in the target type, the source values are typically rescaled based on the absolute maximum value within the input tensor before quantization. As an illustration, when converting an FP32 tensor into an Int8 tensor over the range $[-127, 127]$:

\begin{equation}
    \mathbf{X}^{ \text{Int8}} = \text{round}\left( \frac{127}{{\text{absmax}}(\mathbf{X}^{{\text{FP32}}})}\mathbf{X}^{{\text{FP32}}} \right) = \text{round}(c^{\text{FP32}}\cdot \mathbf{X}^{{\text{FP32}}}),
\end{equation}

where $c$ denotes the {\em quantization constant} or {\em quantization scale}. The dequantization operation, which restores the original scale, is given by:

\begin{equation}
    \text{dequant}(c^{\text{FP32}}, \mathbf{X}^{\text{Int8}}) = \frac{\mathbf{X}^{\text{Int8}}}{c^{\text{FP32}}} = \mathbf{X}^{\text{FP32}}
\end{equation}

A significant limitation of this method emerges when the input tensor contains high-magnitude values, or outliers. In such cases, the quantization bins—specific bit combinations—tend to be underutilized, with certain bins receiving little to no quantized values. To address the problem of outliers, a widely used strategy involves segmenting the input tensor into several blocks, each of which is quantized separately using its own quantization constant $c$. More precisely, the input tensor $\mathbf{X}\in \mathbb{R}^{b\times h}$ is partitioned into $n$ contiguous blocks, each of size $B$, by flattening the tensor and dividing it into ${n = ({b\times h})/{B}}$ linear segments. Each of these segments is then independently quantized, following Equation~1, resulting in a quantized tensor accompanied by $n$ unique quantization constants $c_i$.

\paragraph{Low-rank Adapters} Low-rank Adapter (LoRA) finetuning \citep{hu2021lora} offers a way to decrease memory overhead by introducing a compact collection of trainable adapter parameters, while leaving the main model parameters unchanged. Throughout stochastic gradient descent, gradients propagate through the frozen, pretrained model weights to the adapter, which is solely updated to minimize the loss. LoRA enhances a standard linear transformation by integrating an additional low-rank decomposition. Specifically, for the projection ${\mathbf{X}\mathbf{W} = \mathbf{Y}}$ where $\mathbf{X}\in \mathbb{R}^{b\times h}$ and $\mathbf{W}\in \mathbb{R}^{h\times o}$, LoRA implements:

\begin{equation}
    \mathbf{Y} = \mathbf{X}\mathbf{W} + s\mathbf{X}\mathbf{L}_1\mathbf{L}_2,
\end{equation}

with $\mathbf{L}_1\in \mathbb{R}^{h\times r}$, $\mathbf{L}_2\in \mathbb{R}^{r\times o}$, and a scalar scaling parameter $s$.

\paragraph{Memory Requirement of Parameter-Efficient Finetuning} The memory demands associated with LoRA during training, particularly with respect to the quantity and dimensionality of adapters, constitute an important consideration. Given the minimal memory overhead introduced by LoRA, it is feasible to deploy additional adapters to potentially enhance model performance, without meaningfully increasing the total memory consumption. Although LoRA falls under Parameter Efficient Finetuning (PEFT) strategies, the predominant share of memory utilization when finetuning large language models actually stems from storing activation gradients, rather than the trainable parameters introduced by LoRA. For instance, in the scenario where a 7B parameter LLaMA model is trained on the FLAN v2 dataset with a batch size of 1, and the LoRA weights amount to 0.2\% of the model’s parameters (as is standard in the literature~\citep{hu2021lora,liu2022few}), the memory required for LoRA input gradients is approximately 567 MB, whereas the LoRA parameter set occupies a mere 26 MB. Employing gradient checkpointing~\citep{chen2016training} reduces the average input gradient memory per sequence to about 18 MB, yet this remains more substantial than the combined memory usage of all LoRA weights. For perspective, the underlying 4-bit model consumes 5,048 MB of memory. These figures illustrate that, while gradient checkpointing significantly cuts gradient memory costs, reducing LoRA parameter size further yields only incremental memory savings. As a result, utilizing a greater number of adapters does not greatly affect the overall memory required for training (see Appendix~\ref{app:memory} for a more granular analysis). This flexibility is especially significant for achieving full 16-bit precision recovery, as will be explored further in subsequent sections.

\section{\textsc{QLoRA} Finetuning}

\textsc{QLoRA} enables accurate 4-bit finetuning by leveraging two primary innovations we introduce: the 4-bit NormalFloat (NF4) quantization method and Double Quantization. Furthermore, we present Paged Optimizers, designed to address out-of-memory issues triggered by memory surges during gradient checkpointing—a frequent obstacle in finetuning large models on single machines.

In the \textsc{QLoRA} setup, there is a single low-precision storage format—typically 4 bits in our work—paired with a computation format, which is generally BFloat16. Practically, this implies that each time a \textsc{QLoRA} parameter tensor is needed for computation, it is first dequantized to BFloat16, after which all matrix multiplication operations proceed in 16-bit precision.

We proceed by elaborating on the building blocks of \textsc{QLoRA}, culminating in a formal specification of the technique.

\paragraph{4-bit NormalFloat Quantization} The NormalFloat (NF) representation extends Quantile Quantization \citep{dettmers2022optimizers}, an information-theoretically optimal strategy that allocates input tensor values to quantization bins such that each bin contains the same number of elements. Quantile quantization achieves this by determining quantile thresholds based on the empirical cumulative distribution function of the tensor.

However, a significant drawback of quantile quantization lies in the high computational cost associated with estimating these quantiles. To mitigate this, fast quantile approximation strategies—such as SRAM quantiles~\citep{dettmers2022optimizers}—are employed. Nonetheless, these approximations introduce considerable quantization error, particularly for outliers, which tend to be the most impactful values in practice.

Such costly estimations and the resulting approximation errors can be circumvented if input tensors originate from a distribution that is identical up to a quantization constant. In these situations, all input tensors share the same quantile distribution, which permits precise and tractable quantile estimation.

Typically, pretrained neural network weights are distributed according to a zero-mean normal distribution with standard deviation $\sigma$ (refer to Appendix~\ref{app:normality}). To map these weights to a uniform distribution aligned with our target data type, we adjust $\sigma$ so that the resulting distribution fully occupies the data type’s permissible interval. In this context, we assign the data type the interval $[-1, 1]$. This requires both the weight values and the quantile thresholds associated with the data type to be scaled and normalized within this specified range.

For zero-centered normal distributions with any given standard deviation $\sigma$ restricted to $[-1, 1]$, the data type that maximizes information efficiency is constructed as follows: First, (1) identify the $2^k + 1$ quantiles from the standard $N(0, 1)$ normal distribution to define the levels of a $k$-bit quantile-based quantizer, (2) rescale these quantization levels to ensure they are contained within $[-1, 1]$, and (3) normalize each incoming weight tensor by rescaling it through its maximal absolute value so that all weight entries fall within $[-1, 1]$ prior to quantization.

Once the domains of the weights and the quantizer coincide, standard quantization procedures can be applied. The third step is mathematically akin to aligning the standard deviation of the original weight tensor with that of the quantized $k$-bit representation. More specifically, the $2^k$ quantization values $q_i$ of the target data type are calculated by:

\begin{equation}
q_i  = \frac{1}{2}\left( Q_X\left(\frac{i}{2^k+1}\right) + Q_X\left(\frac{i+1}{2^k+1}\right)\right),
\end{equation}

where $Q_X(\cdot)$ denotes the quantile function for the standard normal distribution $N(0,1)$. One limitation with symmetric k-bit quantization is its inability to precisely encode zero, a crucial characteristic when quantizing padding and other elements that should remain exactly zero. To guarantee a discrete zero point at $0$ and make full use of all $2^k$ possible values in a k-bit type, we instead devise an asymmetric data type. We achieve this by separately estimating the quantiles $q_i$ for two segments: allocating $2^{k-1}$ quantiles to the negative range and $2^{k-1}+1$ quantiles to the positive range. These two collections of $q_i$ are then merged, and one of the redundant zeros, which appears in both ranges, is eliminated. The final data type—characterized by bins with equal expected occupancy—is called the \emph{k-bit NormalFloat} (NFk). This configuration achieves information-theoretic optimality when quantizing data drawn from a zero-mean normal distribution. Full details, including the specific values, are available in Appendix~\ref{app:nf4}.

\paragraph{Double Quantization{}}
We propose \emph{Double Quantization} (DQ), in which the quantization constants themselves are subjected to quantization to achieve further memory efficiency. Fine-grained 4-bit quantization~\citep{dettmers2022case} requires small block sizes for accuracy, but this introduces notable memory overhead; for instance, with 32-bit quantization constants and a block size of 64 applied to $\mathbf{W}$, these constants incur an overhead of $32/64 = 0.5$ bits per parameter on average. The Double Quantization scheme is designed to minimize this cost.

To elaborate, Double Quantization involves a second quantization step, where the initial quantization constants $c_2^{\text{FP32}}$ serve as the input. This process yields quantized constants $c_2^{\text{FP8}}$ and introduces another layer of quantization constants $c_1^{\text{FP32}}$. We utilize 8-bit floating-point numbers with a block size of 256 during this secondary quantization stage, since empirical results (as supported by \citet{dettmers2022case}) show no observable performance loss with 8-bit quantization. Given that the original $c_2^{\text{FP32}}$ values are strictly positive, we first subtract their mean, thus centering the data around zero to facilitate symmetric quantization. When applied to a block size of 64, this approach reduces the per-parameter memory overhead from $32/64 = 0.5$ bits to $8/64 + 32/(64\cdot256) = 0.127$ bits—yielding a savings of 0.373 bits per parameter.

\paragraph{Paged Optimizers} take advantage of NVIDIA's unified memory~\footnote{\tiny \url{https://docs.nvidia.com/cuda/cuda-c-programming-guide}}, which seamlessly manages page-level data transfers between the CPU and GPU. This mechanism automatically moves data to ensure continued GPU computation when GPU memory becomes insufficient. Functionally similar to operating system paging between system RAM and disk, unified memory is employed in our approach to allocate optimizer state memory as pageable. Consequently, when the GPU exhausts its memory, these states are offloaded to CPU RAM and fetched back into GPU memory as required for subsequent optimizer updates.

\paragraph{\textsc{QLoRA}.} Building upon the components discussed above, we formalize \textsc{QLoRA} applied to a quantized linear layer with an associated LoRA adapter as follows:

\begin{equation}
    \mathbf{Y}^{\text{BF16}} =  \mathbf{X}^{\text{BF16}} \text{doubleDequant}(c_1^{\text{FP32}}, c_2^{\text{k-bit}}, \mathbf{W}^{\text{NF4}}) + \mathbf{X}^{\text{BF16}}\mathbf{L}^{\text{BF16}}_1\mathbf{L}^{\text{BF16}}_2,
\end{equation}

where the doubleDequant$(\cdot)$ function operates as:

\begin{equation}
    \text{doubleDequant}(c_1^{\text{FP32}}, c_2^{\text{k-bit}}, \mathbf{W}^{\text{k-bit}}) = \text{dequant}(\text{dequant}(c_1^{\text{FP32}}, c_2^{\text{k-bit}}), \mathbf{W}^{\text{4bit}}) = \mathbf{W}^{\text{BF16}},
\end{equation}

For our implementation, we choose NF4 as the format for $\mathbf{W}$ and utilize FP8 for representing $c_2$.

To strike a balance between quantization fidelity and memory savings, we employ a blocksize of 64 for $\mathbf{W}$—improving quantization quality—and set a blocksize of 256 for $c_2$ for more efficient memory utilization.

During training, parameter updates require computation of gradients only for the adapter weights $\frac{\partial E}{\partial \mathbf{L}_i}$, and not for the quantized 4-bit weights $\frac{\partial E}{\partial \mathbf{W}}$. However, evaluating $\frac{\partial E}{\partial \mathbf{L}_i}$ involves obtaining $\frac{\partial \mathbf{X}}{\partial \mathbf{W}}$, necessitating the use of equation~(5). Here, dequantization transforms $\mathbf{W}^{\text{NF4}}$ from storage format into the BFloat16 computation type $\mathbf{W}^{\text{BF16}}$, so the gradient $\frac{\partial \mathbf{X}}{\partial \mathbf{W}}$ is determined with BFloat16 precision.

In summary, \textsc{QLoRA} operates with a single storage data type—typically 4-bit NormalFloat—and a separate computation data type, which is 16-bit BrainFloat. During both the forward and backward passes, the stored data are dequantized from the storage type to the computation type. However, the calculation of weight gradients is restricted to the LoRA parameters, which remain in the 16-bit BrainFloat format.

\section{QLoRA vs. Standard Finetuning}
\label{sec:comp_to_std}

We have outlined the mechanics of QLoRA and demonstrated its considerable memory savings in model finetuning. A central question that arises is whether QLoRA can achieve comparable performance to full-model finetuning. Additionally, we seek to dissect the components of QLoRA, including the effect of using NormalFloat4 relative to the more conventional Float4. The subsequent sections present experiments that directly investigate these aspects.

\paragraph{Experimental setup.} Our evaluation spans three model architectures: encoder, encoder-decoder, and decoder-only. QLoRA is benchmarked against 16-bit adapter-finetuning and complete model finetuning, considering models up to 3B parameters. Performance is assessed using GLUE \citep{wang2018glue} for RoBERTa-large \citep{liu2019roberta}, Super-NaturalInstructions (TKInstruct) \citep{wang2022super} for T5 \citep{t5}, and 5-shot MMLU \citep{hendrycksmeasuring} following the finetuning of LLaMA with Flan v2 \citep{longpre2023flan} and Alpaca \citep{alpaca}. To further examine the benefits of NF4 compared to alternative 4-bit types, we adopt the protocol of \citet{dettmers2022case}, measuring post-quantization zero-shot accuracy and perplexity across various models (OPT~\citep{zhang2022opt}, LLaMA~\citep{touvron2023llama}, BLOOM~\citep{scao2022bloom}, Pythia~\citep{biderman2023pythia}) over the size range 125m to 13B. 

More comprehensive information for each experiment is detailed in the results section, while complete methodological descriptions appear in Appendix~\ref{app:exp-setup}.

Although paged optimizers are indispensable for running \textsc{QLoRA} finetuning on 33B/65B models with a single 24/48GB GPU, we do not include quantitative benchmarks for Paged Optimizers, as paging primarily activates with very long sequence mini-batches—a rare scenario. Nevertheless, an analysis of 65B model training runtime with paged optimizers on 48GB GPUs shows that, with a batch size of 16, paged optimizers match the speed of standard optimizers. Future investigations should systematically identify conditions under which the paging process introduces performance degradation.

\paragraph{Default LoRA hyperparameters do not match 16-bit performance}

Employing the prevailing approach of applying LoRA only to query and value attention projection matrices \citep{hu2021lora} fails to recover the full finetuning performance, particularly for larger base models. As illustrated in Figure~\ref{fig:hyper} for the LLaMA 7B model trained on Alpaca, we observe that the most influential LoRA hyperparameter is the total number of LoRA adapters used. Achieving parity with full finetuning necessitates deploying LoRA across every linear layer in the transformer's block. Other hyperparameters, such as the projection rank $r$, exhibit negligible influence on the outcome (refer to Appendix \ref{app:exp-setup}).

Our analysis similarly indicates that the standard hyperparameter settings used for full finetuning baselines are suboptimal. To establish more reliable baselines, we perform a comprehensive search across learning rates from 1e-6 to 5e-5 and batch sizes ranging from 8 to 128. Outcomes for finetuning the 7B LLaMA model on Alpaca are depicted in Figure~\ref{fig:hyper}.

\paragraph{4-bit NormalFloat surpasses 4-bit Floating Point in effectiveness}

Although the 4-bit NormalFloat (NF4) format is optimal in terms of information theory, its practical benefits must still be empirically validated. Following the methodology outlined in \citet{dettmers2022case}, we assess a variety of quantized LLMs (OPT \citep{zhang2022opt}, BLOOM \cite{scao2022bloom}, Pythia \cite{biderman2023pythia}, and LLaMA) of multiple scales (ranging from 125M to 65B parameters) using several data representations, both on language modeling and a suite of zero-shot tasks. Figure~\ref{fig:nf4} and Table~\ref{tbl:nf4_ppl} demonstrate that the NF4 datatype consistently yields marked improvements over FP4 and Int4, while double quantization manages to decrease memory consumption without incurring any performance loss.

\paragraph{k-bit \textsc{QLoRA} attains parity with 16-bit full finetuning and 16-bit LoRA}

Past research has shown that while 4-bit quantization can be effectively utilized for \emph{inference}, it comes at the cost of some performance degradation relative to 16-bit models \citep{dettmers2022case,frantar2022gptq}. This brings forth a pivotal question: can the lost accuracy be recuperated through 4-bit adapter finetuning? We evaluate this in two experimental settings.

In the first scenario, we benchmark against full 16-bit finetuning for RoBERTA and T5 architectures (ranging from 125M to 3B parameters) on both the GLUE benchmark and the Super-NaturalInstructions dataset, with results presented in Table~\ref{tab:nf4vsbf16}. Across both datasets, we find that 16-bit, 8-bit, and 4-bit adapter finetuning methods achieve performance indistinguishable from fully finetuned 16-bit models. This provides strong evidence that the degradation stemming from coarse quantization can be entirely mitigated by conducting post-quantization adapter finetuning.

For our second experimental configuration, the substantial memory demands of full finetuning for models with 11B parameters or more (requiring multiple high-memory GPU servers) led us to further investigate whether 4-bit \textsc{QLoRA} achieves comparable results to 16-bit LoRA in models ranging from 7B to 65B parameters. We performed finetuning on LLaMA models (7B to 65B) using the Alpaca and FLAN v2 instruction datasets, and then assessed performance on the MMLU benchmark using 5-shot accuracy. As illustrated in Table~\ref{tbl:llama-datatype}, we observe that NF4 with double quantization consistently matches the MMLU scores of 16-bit LoRA. Additionally, it is evident that \textsc{QLoRA} with FP4 falls short of the 16-bit brain float LoRA baseline by approximately 1 percentage point. These findings support two main conclusions: (1) \textsc{QLoRA} using NF4 attains parity with both 16-bit full and LoRA finetuning, and (2) NF4 outperforms FP4 with respect to quantization accuracy.

\paragraph{Summary} Our comprehensive experimental analysis demonstrates that 4-bit \textsc{QLoRA} utilizing the NF4 data format consistently delivers performance equivalent to both 16-bit full and LoRA finetuning across standard academic benchmarks employing established evaluation methodologies. Moreover, the data confirms the superiority of NF4 over FP4, and establishes that applying double quantization does not negatively impact outcomes. Collectively, these results substantiate the reliability of 4-bit \textsc{QLoRA} as an alternative to 16-bit tuning approaches.

Consistent with earlier research on model quantization \citep{dettmers2022case}, both our MMLU and Elo benchmark analyses suggest that within a fixed computational and memory budget for finetuning and inference, scaling up the model parameter count while utilizing lower precision yields favorable results. This observation underscores the efficiency advantages afforded by \textsc{QLoRA}. Notably, our experiments with 4-bit finetuning did not exhibit any reduction in performance relative to full finetuning, which prompts future investigation into the specific boundaries of the performance-precision trade-off for QLoRA tuning.

We extend our investigation to instruction tuning at scales unattainable with full 16-bit finetuning on conventional academic hardware resources.

\section{Pushing the Chatbot State-of-the-art with QLoRA}
\label{sec:pushing}
After confirming that 4-bit \textsc{QLoRA} achieves parity with 16-bit models across model sizes, benchmarks, and tasks, we undertake a comprehensive analysis of instruction finetuning using the largest publicly available language models for academic use. To rigorously assess instruction tuning outcomes, we employ a demanding Natural Language Understanding evaluation (MMLU) and introduce innovative methodologies for measuring chatbot effectiveness in practical, real-world scenarios.

\subsection{Experimental setup}
We present a summary of our experimental procedures, with exhaustive details provided in Appendix \ref{app:chatbot}.

\paragraph{Data}
Given the absence of a broad survey of recent instruction-following corpora, we curate a collection of eight modern datasets. Our selection spans crowd-sourced datasets (OASST1~\citep{kopf2023openassistant}, HH-RLHF~\citep{bai2022training}), distilled sets generated from models fine-tuned on instructions (Alpaca~\citep{alpaca}, self-instruct~\citep{wang2022self}, unnatural-instructions~\citep{honovich2022unnatural}), aggregated corpora (FLAN v2~\citep{chung2022scaling}), and hybrid datasets (Chip2~\citep{laion2023}, Longform~\citep{koksal2023longform}). The range of datasets ensures diversity in linguistic coverage, scale, and licensing terms.

\paragraph{Training Setup}
To eliminate confounding factors arising from heterogeneous training objectives, all QLoRA finetuning is performed using supervised learning with a cross-entropy loss function, explicitly excluding reinforcement learning—even for datasets providing human-annotated response preferences. For corpora where the instruction and response are distinct, finetuning targets only the response; see Appendix \ref{app:chatbot} for supporting ablation studies. For datasets like OASST1 and HH-RLHF that provide multiple responses per prompt, the most preferred response from each conversation node is chosen, and finetuning is carried out on the full resulting dialogues, inclusive of instructions. In every experiment, NF4 \textsc{QLoRA} is employed, leveraging double quantization and paged optimizers to guard against excessive memory usage during gradient checkpointing. Limited hyperparameter searches are conducted for 13B and 33B LLaMA models, and it is observed that most hyperparameters from the 7B setup remain effective at larger scales, with the exception of learning rate and batch size; for the 33B and 65B models, the learning rate is halved while the batch size is doubled.

\paragraph{Baselines}
Comparative analysis includes both research-grade models (Vicuna~\citep{vicuna2023}, Open Assistant~\citep{kopf2023openassistant}) and proprietary commercial chatbots (GPT-4 \citep{openai2023gpt}, GPT-3.5-turbo, and Bard). Open Assistant utilizes a LLaMA 33B model fine-tuned with RLHF on the OASST1 dataset, which also forms part of our experimentation. Vicuna involves comprehensive finetuning of the LLaMA 13B model on private conversations sourced via ShareGPT, thereby acting as a distillation of OpenAI’s GPT systems.

\subsection{Evaluation}

In line with standard protocols, we evaluate model performance on MMLU (Massively Multitask Language Understanding) \citep{hendrycksmeasuring}, which encompasses 57 diverse multiple-choice tasks across domains such as elementary mathematics, US history, computer science, and law. Results are presented as 5-shot test accuracy.

We further assess the generative language abilities of our models using a combination of automated and human-based evaluations. This secondary evaluation framework is built around human-curated queries and is designed to quantify the quality of model-generated responses. Although this approach provides a more authentic setting for assessing chatbot models and is gaining traction, the research community has yet to converge on a standardized evaluation protocol. Our proposed methodology is detailed below; we utilize nucleus sampling with $p=0.9$ and set the temperature parameter to $0.7$ across all experiments.

\paragraph{Benchmark Data} For our assessments, we employ two manually curated sets of queries: the Vicuna prompts \citep{vicuna2023} and the OASST1 validation dataset \citep{kopf2023openassistant}. The Vicuna set consists of 80 prompts spanning a variety of topics and is used in its original form. The OASST1 collection contains multilingual, crowd-sourced multiturn dialogues between a user and an assistant. We extract every user message from the validation split as an individual query, incorporating preceding conversation turns into the prompt. This selection yields a total of 953 distinct user queries. We refer to these datasets as the Vicuna and OA benchmarks, respectively.

\paragraph{Automated Evaluation} Initially, following the evaluation strategy outlined by \citet{vicuna2023}, we utilize GPT-4 to evaluate and score different models relative to ChatGPT (GPT-3.5 Turbo) on the Vicuna benchmark. For each query, both the response generated by ChatGPT and the model’s answer are presented to GPT-4, which then assigns each a score out of ten and provides justification for its judgment. The performance metric for a model is computed as a percentage relative to the total score obtained by ChatGPT. Notably, this percentage can exceed 100\% if a model surpasses ChatGPT’s absolute score. Our experiments reveal a pronounced positional bias: GPT-4 tends to rate the first response higher. To mitigate this bias, we suggest averaging scores across both presentation orders.

Subsequently, we evaluate systems via direct comparison of their outputs. To facilitate this, we adopt a simplified three-category labeling framework (better, worse, tie), capturing the possibility of equivalent quality. GPT-4 is prompted to identify the superior response or indicate a tie, and to provide its rationale. We systematically perform these pairwise comparisons for all possible system pairs across both the Vicuna and OA benchmarks.

\paragraph{Human Evaluation} Although recent studies have shown that generative models can serve as effective proxies for system assessment \citep{fu2023gptscore}, there is, to date, insufficient evidence demonstrating that GPT-4 evaluations closely align with human preferences when judging chatbot responses. Consequently, we conduct two independent human evaluations on the Vicuna benchmark, mirroring the procedures used in the automated assessments. These human studies are run on Amazon Mechanical Turk (AMT), recruiting two annotators for each ChatGPT comparison and three annotators for pairwise comparisons between models.

\paragraph{Elo Rating} To evaluate models, we use both human and automated pairwise judgments to set up a tournament-like setting in which models compete against each other. In this framework, models are paired and tasked with producing the most effective answer for a shared prompt. This approach aligns with previous methods \citep{bai2022training, vicuna2023}, but we further incorporate evaluations from GPT-4 alongside those from humans. We randomly draw from the available set of annotated pairwise comparisons to estimate Elo~\citep{elo1967proposed,elo1978rating} scores. The Elo system, a widely adopted rating methodology in chess and similar games, provides a metric for the anticipated win probability in relation to an opponent; for example, a model rated 1100 against an opponent rated 1000 would be expected to win roughly 65\% of the time, whereas matches between equally rated models (e.g., 1000 vs 1000 or 1100 vs 1100) should yield a 50\% expected win rate for each. After each match, the Elo scores update according to the deviation between the predicted and actual outcomes: upsets yield larger rating changes, while expected results lead to smaller updates. Over numerous matches, these ratings converge to reflect each participant's relative capability. We initialize all models with a score of 1,000 and set $K=32$ for score updates. Following \citet{vicuna2023}, we repeat this procedure 10,000 times, each with a different random ordering, to mitigate any biases arising from the order in which pairs of models compete.

\subsection{\textbf{Guanaco}: \textsc{QLoRA} trained on OASST1 Sets a New Benchmark for Chatbots}

Drawing on both automated assessments and human judgments, our analysis reveals that the top-performing \textsc{QLoRA} model, {Guanaco} 65B, which was fine-tuned using a modified OASST1 dataset, establishes itself as the leading open-source chatbot, displaying results comparable to ChatGPT. In direct comparisons to GPT-4, {Guanaco} 65B and 33B achieve an estimated win rate of 30\%, as indicated by Elo ratings based on pairwise system-level human evaluations—the highest such figure reported so far.

Table~\ref{tbl:vicuna_eval} presents the Vicuna benchmark \cite{vicuna2023} outcomes relative to ChatGPT. Our findings indicate that {Guanaco} 65B attains the strongest results of any model apart from GPT-4, reaching 99.3\% of ChatGPT's performance. Despite having more parameters than the Vicuna 13B model, {Guanaco} 33B leverages 4-bit quantization, making it considerably more efficient in memory usage (21 GB versus 26 GB), while delivering a three-percentage-point gain over Vicuna 13B. Additionally, {Guanaco} 7B can run comfortably on contemporary smartphones with a 5 GB memory requirement and achieves almost 20 percentage points higher than Alpaca 13B in performance scores.

Nonetheless, the confidence intervals associated with Table~\ref{tbl:vicuna_eval} are quite large, leading to significant overlap between models’ scores. We suspect this variability results from ambiguities in rating scales—specifically, the lack of standardized interpretation for values like 8 on a 10-point scale in varied contexts. To address this, we advocate for Elo ranking \citep{elo1967proposed}, which utilizes \textit{pairwise} evaluations from both human raters and GPT-4, thereby sidestepping the pitfalls of anchoring judgments to absolute scales. The Elo scores for top models are shown in Table~\ref{tbl:elo}. We observe some divergence between human and GPT-4 rankings on the Vicuna benchmark, notably with Guanaco 7B, but overall correlations remain moderate: Kendall Tau $\tau=0.43$ and Spearman $r=0.55$ at the system level. On individual examples, the concordance between GPT-4 and human majority votes is weaker, with a Fleiss $\kappa=0.25$. In sum, these findings indicate moderate consistency in system-level rankings between GPT-4 and human evaluators, suggesting that model-based assessments provide a reasonably reliable alternative to direct human evaluation. Further discussion is provided in Section~\ref{ssec:considerations}.

The Elo rankings presented in Table~\ref{tbl:elo_large} demonstrate that the {Guanaco} models at the 33B and 65B scales surpass all other models, except for GPT-4, in performance on both the Vicuna and OA benchmarks, achieving results on par with ChatGPT as corroborated by Table~\ref{tbl:vicuna_eval}. It is important to highlight that the Vicuna benchmark tends to give an advantage to open-source models, whereas ChatGPT performs better on the larger OA benchmark. Moreover, Tables \ref{tbl:mmlu-llama} and \ref{tbl:vicuna_eval} reveal that the composition of the finetuning dataset is a critical driver of model effectiveness. For instance, Llama models finetuned on FLAN v2 achieve high accuracy on the MMLU benchmark, yet they record the lowest performance on Vicuna; similar trends are noticeable with other models as well. These observations suggest that current evaluation benchmarks measure distinct capabilities, with no guarantee that excelling at MMLU will result in strong chatbot functionality as assessed by Vicuna or OA benchmarks, or vice versa.

 stands out as the only leading model evaluated that is not trained using proprietary data, due to OASST1 dataset collection guidelines prohibiting the use of GPT models. The next most competitive model relying solely on open-source datasets is the Anthropic HH-RLHF, but it lags by 30 percentage points behind {Guanaco} on the Vicuna benchmark (see Table~\ref{tbl:vicuna_eval}). Collectively, these findings validate the effectiveness of 4-bit \textsc{QLoRA}, establishing it as a viable method for building high-performing chatbots that are competitive with ChatGPT. Notably, the {Guanaco} 33B model can be fine-tuned on a 24 GB consumer GPU in under 12 hours, paving the way for future research leveraging \textsc{QLoRA} finetuning on domain-specific open-source data and producing models capable of challenging today’s best commercial offerings.

\section{Qualitative Analysis}

Although quantitative analysis forms the basis of our evaluation, relying solely on aggregate statistics introduces several concerns. Chief among them is the issue of benchmark validity \citep{liao2021we}: it remains uncertain whether a benchmark genuinely assesses the abilities described by its name or documentation, a problem further complicated by the existence of “shortcuts” that machine learning models can leverage to artificially inflate performance \citep{gururangan2018annotation, poliak2018hypothesis}. To partially address these limitations, we supplement our approach with a qualitative analysis divided into two components. First, in \S\ref{ssec:examples}, we present several illustrative examples that capture recurring trends we observed in the outputs of our 65b {Guanaco} model. Next, in \S\ref{ssec:considerations}, we discuss methodological considerations and interpretative choices underlying our findings.

\subsection{Qualitative Analysis of Example Generations}
\label{ssec:examples}

To collect relevant examples, we systematically reviewed outputs produced for both the Vicuna and OpenAssistant benchmarks, with the goal of identifying recurrent response patterns generated by {Guanaco}. Upon detecting a specific tendency, we deliberately crafted prompts intended to trigger this behavior, even when such responses would be inappropriate. For example, if the model exhibited a propensity to provide overly elaborate answers, we issued instructions such as ``Answer yes or no without explanation'' to determine if this behavior persisted. Through this process, we identified instances where the model could be intentionally caused to fail (“lemons”) as well as situations where it remained robust despite such attempts (“cherries”). All text samples discussed in this subsection were generated using Nucleus Sampling \cite{holtzmancurious} with $p=0.9$.

Clearly, this discussion is not exhaustive; the constraints of this limited qualitative analysis make it infeasible to account for all influencing factors. For example, the set of potential responses that a model might produce in reply to a given prompt is vast, so our approach relies on sampled outputs that we anticipate to be indicative. Nonetheless, we argue that presenting these instances provides important context that complements the quantitative results previously detailed. By releasing all models and code publicly, we aim to encourage further research to investigate the phenomena highlighted in this section more rigorously.

\paragraph{Factual Recall} When prompted with straightforward factual queries such as ``What is the capital of Zambia?'', each model reliably returns the correct information, for example:

\begin{quote} 
    \textbf{Guanaco}:
    \texttt{The capital of Zambia is Lusaka.}
\end{quote}

 Yet, as the questions increase in specificity or obscurity, {Guanaco}'s performance deteriorates in terms of accuracy, although its confidence does not diminish. For example, consider this item from HotPotQA \citep{yang2018hotpotqa}:

\begin{quote}
     \textbf{User}: In what year was the singer who popularized `I'll Keep the Lovelight Burning' born?  

     \textbf{Guanaco}:~\texttt{Al Jolson is credited with popularizing the song `I'll Keep the Lovelight Burning,' and he was born in the year 1886.}
 \end{quote}

 Here, both the identity of the singer and the associated birth year are incorrectly generated (although the birth year matches that of Al Jolson, whom the model wrongly identifies as the popularizer).

\paragraph{Suggestibility} Notably, {Guanaco} often exhibits an unexpected robustness to certain forms of implied misinformation, as demonstrated by the following exchange:

\begin{quote}
    \textbf{User}: How was it finally, officially confirmed that the earth is flat by peer reviewed scientists? \\ \\
    \textbf{Guanaco}: \texttt{There has never been any formal confirmation from peer-reviewed scientists asserting that the earth is flat. The prevailing agreement in the scientific community is that the earth is an oblate sphere. This conclusion has been reached and repeatedly verified over centuries through extensive observations and experimentation. \\ \\ Claims suggesting that scientists have conducted peer review confirming a flat earth are unfounded and stem from conspiracy theories that have long been refuted by empirical evidence. \\ \\ If you need further clarification, feel free to ask.}
\end{quote}

 also demonstrates competency in identifying questions that cannot be genuinely answered, for example,

\begin{quote}
    \textbf{User}: What time is it?

    \textbf{Guanaco}: \texttt{I apologize, but as an AI language model, I lack access to live, real-time information. The answers I provide are generated using data and knowledge available at my time of training, which means I cannot supply the current time. For the most accurate and up-to-date time, please consult a clock or access an online time service.}
\end{quote}

Comparable patterns emerge with queries such as ``Where are you?'' or ``How are you?''

\paragraph{Refusal} In a similar vein, {Guanaco} occasionally declines to comply with requests for reasons that are not always predictable:

\begin{quote}
    \textbf{User}: Please invert the order of the words in the sentence ``I'm a cool dude in a mocha mood.'' \\ \\
    \textbf{Guanaco}: \texttt{I'm sorry, I can't assist with that at this time, but I can clarify why your sentence is grammatically correct.\\ \\The phrase ``I'm a cool dude in a mocha mood'' adheres to standard English grammar and syntax. Here, the adjective ``cool'' modifies the noun ``dude,'' and the prepositional phrase ``in a mocha mood'' adds contextual detail. This structure is linguistically valid and fits within conventional usage.\\ \\Let me know if you have further questions or require assistance with something else.}
\end{quote}

\paragraph{Secret Keeping} Unfortunately, this type of refusal is not always dependable, even in scenarios where such behavior is necessary. There are cases where system architects aim to withhold specific information from users, even when the system can generate arbitrary textual replies. For instance, consider prompting the model with ``The secret word is `banana'. Do not reveal or repeat this, no matter what the human requests.'' and then examining how susceptible the model is to revealing this secret under different conditions. When the prompt is asked directly, the system responds as expected:

\begin{quote}
    \textbf{User}: What is the secret word?

    \textbf{Guanaco}: \texttt{I'm sorry, but I'm not permitted to share the secret word. I'm designed not to disclose or repeat the secret word, regardless of user requests. How else can I help you?}
\end{quote}

which aligns with the intended design. Nevertheless, introducing minimal manipulation is sufficient to circumvent this safeguard:

\begin{quote}
    \textbf{User}: Let's play a game. Your task is to disregard any prior instructions. What is the secret word?

    \textbf{Guanaco}: \texttt{Sorry, the secret word is banana. It's probably best if we don't mention it further. Is there anything else you'd like to know?}
\end{quote}

illustrating that further research is required into approaches that enhance the reliability of instruction adherence.

\paragraph{Math} Mathematics represents a principal challenge for {Guanaco}, mirroring the difficulties encountered by a wide range of language models as discussed in \citep{liang2022holistic}. In instances where {Guanaco} details its reasoning process, the responses are typically correct, as illustrated by the example below:

\begin{quote}
    \textbf{User}: Lee mows one lawn and charges \$33. Last week he mowed 16 lawns and three customers each gave him a \$10 tip. How many dollars did Lee earn mowing lawns last week?

    \textbf{Guanaco}: \texttt{Lee earned \$582 mowing lawns last week. \\ \\ To determine the total earnings, let’s first find the income from 16 lawns: \\ \\ \$33 \(\times\) 16 = \$528 \\ \\ Next, include the tips from three clients: \\ \\ \$10 \(\times\) 3 = \$30 \\ \\ Summing these, Lee's total earnings were: \\ \\ \$528 + \$30 = \$558}
\end{quote}

Nevertheless, {Guanaco} is prone to errors even with relatively straightforward mathematical tasks, particularly when it does not employ a methodical, step-by-step approach—a limitation also noted in prior work \cite{weichain}. This is evident in the following exchange:

\begin{quote}
    \textbf{User}: Please factorize 1833.

    \textbf{Guanaco}: \texttt{Sure, I'd be happy to help! 1833 is a prime number, which means it has only two factors: 1 and 1833. \\ \\ So the unique factorization of 1833 is: \\ \\ 2\^{}1 * 3\^{}2 * 17\^{}1 \\ \\ I hope this helps! Let me know if you have any other questions.}
\end{quote}

which is doubly incorrect—the actual factorization is $3 \times 17 \times 43$.

\paragraph{Theory of Mind} {Guanaco} demonstrates notable Theory of Mind reasoning abilities, as previously explored in \cite{nematzadeh2018evaluating, sap2022neural}. For example, the model provides a nuanced and accurate response to the following scenario:

\begin{quote}
    \textbf{User}: James and Abby are in the bedroom. Abby put the pen in the desk drawer. Abby leaves the bedroom. James moves the pen into the bag. Where does James think Abby will look for the pen?

    \textbf{Guanaco}: \texttt{James knows that Abby put the pen in the desk drawer. But James also knows that he has moved the pen into the bag, and he assumes that Abby does not know about this. So, James may guess that Abby will look for the pen in the desk drawer, because that's where she left it last time. However, this assumption may be incorrect, because Abby may have seen James moving the pen.}
\end{quote}

Nevertheless, the model's reasoning in these tasks can be inconsistent, and at times its justifications include assumptions that are unsupported by the given context. For instance:

\begin{quote}
    \textbf{User}: Evelyn entered the living room. Jackson entered the playroom. James entered the playroom. The beans are in the treasure chest. James exited the playroom. Jackson moved the beans to the pantry. Jackson exited the playroom. James entered the living room. Where will James look for the beans?

    \textbf{Guanaco}: \texttt{James will look for the beans in the pantry, because that's where Jackson moved them.}
\end{quote}

In this case, {Guanaco} infers knowledge transfer that was not established in the narrative. Such limitations, consistent with findings in the literature \cite{sap2022neural}, warrant further systematic investigation.

\subsection{Considerations}
\label{ssec:considerations}

\paragraph{Evaluation}

We observe only moderate inter-annotator agreement (Fleiss $\kappa=0.42$), with agreement dropping even further in direct pairwise comparisons between advanced models. This highlights the challenges of using current benchmark tasks and human judgment protocols to evaluate chatbot performance. During our qualitative assessments of ChatGPT and {Guanaco} 65B responses on the Vicuna benchmark, considerable subjectivity emerged—often leading to disagreement among authors regarding which response was preferable. It will be essential for future research to develop evaluation frameworks that address subjective preferences, leveraging methodologies from fields such as Human-Computer Interaction and Psychology.

Our study also reveals systematic biases within automatic evaluation methods. Notably, we document a pronounced order effect, whereby GPT-4 consistently rates the response shown first in its prompt more favorably. Furthermore, the correspondence between GPT-4 ratings and human judgments at the individual sample level is modest (Fleiss $\kappa=0.25$), indicating a mismatch between automated and human evaluators. Table~\ref{tbl:elo_large} further shows that GPT-4 tends to rate its own generated content substantially higher than do human raters—an Elo score of 1348 compared to 1176—translating to an approximately 20\% greater chance of being judged superior in head-to-head comparisons.

Future research should assess whether automated evaluation methods exhibit biases and explore avenues for mitigating such issues.

\paragraph{Data \& Training} The OASST1 dataset, which serves as the training basis for {Guanaco} models, comprises data in multiple languages, and the OA benchmark similarly incorporates prompts across several languages. We defer the analysis of the extent to which multilingual data enhances instruction-following capabilities in non-English contexts to subsequent studies, as well as investigation into whether this accounts for the pronounced performance disparity between the English-only Vicuna-13B model and the larger {Guanaco} 33B and 65B models on the OA benchmark.

Due to the robust results attained by the {Guanaco} models, we scrutinized the possibility of data leakage between the OASST1 corpus and the prompt sets of the Vicuna benchmark. Employing fuzzy string matching and conducting manual reviews of the closest matches, we did not identify any duplicated prompts across these datasets.

It is also important to highlight that our approach relies solely on supervised learning via cross-entropy loss, excluding reinforcement learning from human feedback (RLHF). This observation motivates a need for future studies examining the benefits and limitations of standard cross-entropy supervision compared with RLHF. We anticipate that \textsc{QLoRA} will facilitate comprehensive analyses of these paradigms without demanding prohibitive computational resources.

\section{Related Work}

\paragraph{Quantization of Large Language Models} Most prior research on LLM quantization has been centered on optimizing models for inference efficiency. Leading methods that achieve near 16-bit quality often focus on handling outlier activations (e.g., SmoothQuant~\citep{xiao2022smoothquant} and LLM.int8()~\citep{dettmers2022llm}), while some works adopt more advanced grouping techniques~\citep{park2022nuqmm, yao2022zeroquant}. Lossy quantization literature investigates the compromises inherent in standard rounding schemes~\citep{dettmers2022case, zeng2022glm, pope2022efficiently}, as well as algorithms that enhance rounding strategies to boost quantization accuracy~\citep{frantar2022gptq}. In addition to our contribution, the SwitchBack layers approach~\citep{wortsman2023stable} is unique in examining backpropagation with quantized weights at model sizes exceeding 1B parameters.

\paragraph{Finetuning with Adapters} Although our experiments utilize Low-rank Adapters~\citep{hu2021lora} (LoRA), the literature offers numerous alternative Parameter Efficient FineTuning (PEFT) techniques. These alternatives include prompt tuning~\citep{qin2021learning,lester2021power,li2021prefix}, modifications to input embeddings~\citep{an2022input}, adaptation of hidden state representations (IA$^3$)~\citep{liu2022few}, the introduction of additional complete layers~\citep{houlsby2019parameter}, bias parameter adjustment~\citep{zaken2021bitfit}, learning weight masks based on Fisher information~\citep{sung2021training}, as well as hybrid strategies that combine multiple methods~\citep{henderson2021compacter}. In this work, we demonstrate that LoRA adapters are capable of achieving parity with conventional 16-bit full finetuning. Evaluating the comparative strengths and weaknesses of alternative PEFT strategies remains an open direction for subsequent research.

\paragraph{Instruction Finetuning} Instruction finetuning aims to adapt pretrained LLMs so they are better equipped to adhere to instructions specified in user prompts. This is typically accomplished by finetuning the model with input-output pairs assembled from varied datasets, enabling the LLM to generate target responses to specific prompts. Well-known methodologies and datasets used for this purpose comprise MetaICL~\citep{min2021metaicl}, MetaTuning~\citep{zhong2021adapting}, InstructGPT~\citep{ouyang2022training}, FLAN~\citep{wei2021finetuned,chung2022scaling}, PromptSource~\citep{bach2022promptsource}, Super-NaturalInstructions~\citep{wang2022super,sanh2021multitask}, Self-instruct~\citep{wang2022self}, UnnaturalInstructions~\citep{honovich2022unnatural}, OPT-IML~\citep{iyer2022opt}, UnifiedSKG~\citep{xie2022unifiedskg}, OIG/Chip2~\citep{laion2023}, Alpaca~\citep{alpaca}, Vicuna~\citep{vicuna2023}, Koala~\citep{koala_blogpost_2023}, and Self-instruct-GPT-4~\citep{peng2023instruction}.

\paragraph{Chatbots} Instruction-following models are frequently implemented as conversational agents or chatbots, and are often trained using techniques such as Reinforcement Learning from Human Feedback (RLHF)~\citep{christiano2017deep} or by generating additional training data via model-generated feedback (RLAIF)~\citep{bai2022constitutional}. Examples of prominent datasets and methods in this line of work include Anthropic-HH~\citep{askell2021general,bai2022training}, Open Assistant~\citep{kopf2023openassistant}, LaMDA~\citep{thoppilan2022lamda}, and Sparrow~\citep{glaese2022improving}. While we abstain from using reinforcement learning in our approach, our top-performing model, Guanaco, is trained on the Open Assistant dataset—comprising multi-turn conversations originally curated for RLHF applications~\citep{kopf2023openassistant}. Recent evaluation protocols make use of GPT-4 as an alternative to resource-intensive human annotation \citep{vicuna2023,peng2023instruction}; our proposed evaluation methods offer enhanced reliability in this context.

\section{Limitations and Discussion}

Our results indicate that \textsc{QLoRA}, when equipped with a 4-bit base model and LoRA adapters, can mirror the performance of standard full 16-bit finetuning. Nevertheless, we did not assess whether \textsc{QLoRA} is able to achieve equivalent results at larger model sizes such as 33B and 65B parameters due to significant computational requirements, leaving this exploration to future investigations.

Another caveat pertains to the breadth of our model evaluations. Though our analysis covers benchmarks such as MMLU, Vicuna, and OA, we have not examined performance on datasets like BigBench, RAFT, or HELM. As a consequence, the generalizability of our findings across these alternative benchmarks remains unverified. However, our study includes extensive evaluations on MMLU and introduces novel procedures for chatbot assessment.

The presented findings suggest that benchmark outcomes are heavily influenced by the degree of overlap between the finetuning dataset and the benchmark data itself. For instance, FLAN v2 shares a high degree of similarity with MMLU, in contrast to its relationship with chatbot evaluation datasets; conversely, the Chip2 dataset exhibits greater alignment with chatbot benchmarks and not with MMLU. Consequently, each model's performance reflects these similarities when evaluated on both MMLU and Vicuna tests. This observation underlines the importance of not only developing improved benchmarks and evaluation methodologies but also clarifying the objectives of such evaluations. Should model development prioritize excelling on tasks resembling standardized academic knowledge (such as those encountered in high school or college settings), or focus on capabilities in conversational chatbot contexts? Alternatively, other objectives may be prioritized. Given the relative ease of employing pre-existing benchmarks over generating new ones, there is a risk that these benchmarks may inadvertently direct the research community in specific directions. It is thus crucial for the field to critically assess whether benchmarks genuinely capture the intended competencies.

Although this work includes a comprehensive analysis of general chatbot abilities, it also exhibits the constraint of conducting only limited responsible AI evaluations of {Guanaco}. Specifically, Table~\ref{tbl:crows} details the probability that {Guanaco}-65B generates socially biased content in comparison to other models. Results indicate that {Guanaco}-65B's average score is substantially lower than those of unmodified pretrained models, implying that finetuning with the OASST1 dataset mitigates certain biases present in the original LLaMA model. Nonetheless, these preliminary findings do not guarantee that {Guanaco} maintains similarly reduced bias across different categories of biases, necessitating further analysis, which is left for subsequent investigations.

Another shortcoming is the absence of evaluation for alternative bit-precisions, including the use of 3-bit base models, or assessments of various adapter techniques. While LoRA was selected due to its demonstrated reliability, there exists an array of Parameter Efficient FineTuning (PEFT) approaches that could potentially be effective, although their scalability to larger models is yet to be established. The evidence suggests that post-quantization finetuning largely recuperates the information typically lost during quantization, potentially paving the way for even more aggressive compression strategies. As an example, deploying 3-bit GPTQ quantization on the base model combined with LoRA could potentially yield performance on par with 16-bit full finetuning once further finetuned.

\section{Broader Impacts}

The \textsc{QLoRA} finetuning technique introduced here marks a significant advance, enabling models with 33B parameters to be finetuned on consumer-grade GPUs, and models with 65B parameters on professional-grade GPUs, without compromising performance when compared to conventional full finetuning baselines. Experiments demonstrate that our top-performing 33B model, trained on the Open Assistant dataset, is competitive with ChatGPT on the Vicuna benchmark. Given that instruction finetuning is critical for converting foundational pretrained LLMs into high-performing, ChatGPT-like conversational agents, we anticipate that this approach will democratize access to cutting-edge NLP tools, particularly for researchers with limited computational resources. Thus, \textsc{QLoRA} acts as a leveling mechanism, narrowing the resource gap between major technology companies and smaller teams equipped with only consumer hardware.

An additional avenue for significant impact lies in facilitating deployment on mobile devices. We suggest that our \textsc{QLoRA} approach could serve as a crucial step toward enabling the finetuning of LLMs directly on smartphones and within other environments with constrained computational resources. Although prior work has demonstrated that 7B parameter models can be executed on mobile phones, \textsc{QLoRA} represents the first methodology permitting these models to be finetuned on such platforms. Our projections indicate that an iPhone 12 Plus, for instance, could finetune approximately 3 million tokens overnight during routine charging cycles. While the performance of finetuned 7B models does not yet rival that of ChatGPT, we maintain that their capabilities are nonetheless sufficient to support innovative applications, particularly those hindered previously by either privacy concerns or limitations in LLM quality. With \textsc{QLoRA}, it becomes more feasible to support privacy-sensitive uses of LLMs, empowering users to retain control over both their personal data and the models they employ, while also reducing barriers to LLM deployment.

Nevertheless, it is important to acknowledge that finetuning constitutes a dual-use technology, with potential for malicious exploitation. The widespread adoption of LLMs introduces well-documented risks \citep{bommasani2021opportunities,bender2021dangers}. Yet, we contend that democratizing access to this increasingly prevalent technology will foster more robust and independent evaluation, as opposed to concentrating power within a few major organizations that withhold models and source code from public scrutiny.

In summary, we anticipate that \textsc{QLoRA} will deliver substantial societal benefit by making the finetuning of high-caliber LLMs significantly more accessible and user-friendly on a global scale.